\section{Pricing American Options}

\subsection{Backward Induction And LSM}

\subsubsection*{American exercise problem}

Simulation generates paths forward, but an American option needs a backward exercise rule. For a put with exercise dates $t_1<\cdots<t_n=T$ and $t_{j+1}-t_j=\Delta t$, simulate paths $\{S(t_1),\ldots,S(t_n)\}$ and define $Y_j$ as the cash flow from $t_j$ onward under the selected policy.

At maturity,
$Y_n=\max(K-S(t_n),0)$.

For earlier dates, the generic cash-flow recursion is

$Y_j=\begin{cases}K-S(t_j),&\text{exercise},\\ e^{-r\Delta t}Y_{j+1},&\text{continue}.\end{cases}$

There is no exercise right before $t_1$, so the American put is priced as a claim paying $Y_1$ at $t_1$:
$P_A(0,S)=\hat{\mathbb{E}}\left[e^{-rt_1}Y_1\mid S_0=S\right],\, \hat P_A(0,S)=\frac1N\sum_{i=1}^Ne^{-rt_1}Y_1^{(i)}$.

\subsubsection*{Continuation value}

At $t_j$, the decision must use information available at $t_j$. Compare immediate exercise $K-S(t_j)$ with the conditional discounted continuation value

$f_j(S(t_j))=\hat{\mathbb{E}}\left[e^{-r\Delta t}Y_{j+1}\mid S(t_j)\right]$.

Thus
$Y_j=\begin{cases}K-S(t_j),&K-S(t_j)\ge f_j(S(t_j)),\\ e^{-r\Delta t}Y_{j+1},&K-S(t_j)<f_j(S(t_j)).\end{cases}$

For out-of-the-money put paths, $K-S(t_j)<0$, set $f_j(S(t_j))=0$ and continue: $Y_j=e^{-r\Delta t}Y_{j+1}$.

\subsubsection*{Least-squares approximation}

Longstaff-Schwartz estimates $f_j$ from simulated paths by least squares. Use only in-the-money paths in the regression, because only they are close to the exercise decision. Approximate
$f_j(S(t_j))\approx\sum_{k=0}^m a_{j,k}[S(t_j)]^k$.

Regress the realized discounted continuation value $e^{-r\Delta t}Y_{j+1}$ on basis functions of $S(t_j)$. Degree-$3$ polynomials are often reasonable; Hermite, Laguerre, and Legendre bases are common alternatives, and prices are usually fairly robust to the basis choice.

\subsection{Three-Date American Put Example}

\subsubsection*{Setup}

Let $S(t)$ follow risk-neutral GBM with $S(0)=10$, $r=3\%$, dividend yield $\delta=0$, $\sigma=40\%$, strike $K=12$, maturity $T=1$, and exercise dates $1/3,2/3,1$. Simulate with
$S(t+\Delta t)=S(t)\exp\left[\left(r-\frac{\sigma^2}{2}\right)\Delta t+\sigma\Delta W(t)\right]$.

Using the 8 sample paths in the lecture, the terminal payoff $Y_3=\max(K-S(1),0)$ gives the European put estimate $2.4343$, a lower bound for the American put.

\subsubsection*{Step at $t=2/3$}

For in-the-money paths, regress $Y_3e^{-r\Delta t}$ on $S(2/3)$ using a quadratic:

$Y_3e^{-r\Delta t}=\hat a_0+\hat a_1S(2/3)+\hat a_2[S(2/3)]^2+\epsilon$.

The fitted continuation value is

$f_2(S)=-82.5347+17.7788S(2/3)-0.9063[S(2/3)]^2$.

Then
$Y_2=\begin{cases}K-S(2/3),&K-S(2/3)\ge f_2(S(2/3)),\\ e^{-r\Delta t}Y_3,&\text{otherwise}.\end{cases}$

Out-of-the-money paths continue, so $Y_2=e^{-r\Delta t}Y_3$ for those paths.

\subsubsection*{Step at $t=1/3$}

Repeat the regression using $Y_2e^{-r\Delta t}$ and $S(1/3)$. The fitted continuation value is
$f_1(S)=-8.9488+3.3104S(1/3)-0.2036[S(1/3)]^2$.

Use
$Y_1=\begin{cases}K-S(1/3),&K-S(1/3)\ge f_1(S(1/3)),\\ e^{-r\Delta t}Y_2,&\text{otherwise}.\end{cases}$

Finally estimate the current price by averaging $e^{-r\Delta t}Y_1$. In the lecture example, $\hat P_A=3.0919$, which is above the European estimate $2.4343$.
